% ---- pg-524 （书页 512） ----

(e) 假定 $p_1=0.05,p_2=0.95$ 及 $c(1|2)=c(2|1)$, 重复(c)和(d). 这个先验概率假定是否合理? 请加以解释.

(f) 用(b)中的结果构造联合协方差矩阵 $\mathbf{S}_p$, 并构造\hyperref[eq:11-19]{式(11-19)}中的费希尔线性判别函数. 用这个函数对样本观测值进行分类并计算 APER. 在这种情况下, 费希尔的线性判别函数是否是分类量的一种明智选择? 请解释.

(g) 利用观测值对 $(x_1,x_3)$ 及 $(x_1,x_4)$ 重复(b)到(e). 是否某些变量作为分类量看来比其他变量更好? 请解释.

(h) 对所有四个变量 $(X_1,X_2,X_3,X_4)$ 重复(b)到(e).

\noindent 11.25\quad 约翰逊(Johnson, 见文献[19])曾对\hyperref[tab:11.4]{表 11.4} 中的年度财务数据进行了分析, 旨在发现在判别分析中有影响的观测值. 考虑变量 $X_1=\mathrm{CF}/\mathrm{TD}$ 和 $X_3=\mathrm{CA}/\mathrm{CL}$.

\noindent (a)\quad 利用变量 $X_1$ 和 $X_3$ 的数据构造费希尔的线性判别函数. 用这个函数对样本观测值进行分类并计算 APER[见\hyperref[eq:11-35]{式(11-35)}和\hyperref[eq:11-30]{式(11-30)}]. 在 $(x_1,x_3)$ 坐标系中对数据和判别线作图.

\noindent (b)\quad 约翰逊认为, 第 16 行中的破产企业数据与第 13 行中的健康企业数据是有影响的多元观测值. 用 $X_1$ 和 $X_3$ 的数据, 但将第 16 个数据删除, 计算费希尔线性判别函数. 重复上述计算, 但删除第 13 个数据. 将两条判别线都画在(a)中得出的数据散布图上, 并计算 APER, 计算时略去被删除的点. 删除这两个多元观测值中的任一个是否会造成差别? (注意, 这两个有潜在影响的数据点中的任一个离各自所在的数据群中心都不特别“远”.)

\noindent 11.26\quad 用\hyperref[tab:11.4]{表 11.4} 中的数据定义一个二值响应变量 $Z$: 对破产企业 $Z=0$, 对非破产企业 $Z=1$. 令 $X=\mathrm{CA}/\mathrm{CL}$, 考虑 $Z$ 对 $X$ 的直线回归. 虽然二值响应变量并不满足标准的回归模型假定, 用一个或多个预测变量的线性函数来预测响应变量仍是合理的.

\noindent (a)\quad 用最小二乘法来确定 $Z,X$ 数据的拟合直线. 将破产企业的拟合值作为 $[0,1]$ 区间的点作图. 对非破产企业重复同样的工作, 并将两个图重叠. 一个合理的判别法则是: 若 $Z$ 变量的拟合值接近 0(即拟合值小于 0.5), 则预测一个企业将破产; 类似的, 若 $Z$ 变量的拟合值大于 0.5, 则预测一个企业是健康的. 用这个判决法则对样本企业进行分类. 计算 APER.

\noindent (b)\quad 用所有四个变量 $X_1,\cdots,X_4$(作为预测变量)对 $Z$ 重复以上的线性回归及其他有关的计算. APER 是否有什么改变? 第 16 个观测点(破产企业)与第 13 个观测点(非破产企业)是否还能看成是有影响的?

\noindent (c)\quad 用所有四个变量进行逻辑斯蒂回归.

\noindent 11.27\quad \hyperref[tab:11.5]{表 11.5} 中的数据包含三种鸢尾的 $X_2=$萼片宽度与 $X_4=$花瓣宽度的观测值. 对每种鸢尾有 $n_1=n_2=n_3=50$ 个观测值.

\begin{mvtable}
\phantomsection\label{tab:11.5}\mvtabcap{表 11.5\quad 鸢尾数据}
\mvfit{%
\begin{tabular}{cccccccccccc}
\toprule
\multicolumn{4}{c}{$\pi_1$} & \multicolumn{4}{c}{$\pi_2$} & \multicolumn{4}{c}{$\pi_3$}\\
\cmidrule(lr){1-4}\cmidrule(lr){5-8}\cmidrule(lr){9-12}
萼片长 & 萼片宽 & 花瓣长 & 花瓣宽 & 萼片长 & 萼片宽 & 花瓣长 & 花瓣宽 & 萼片长 & 萼片宽 & 花瓣长 & 花瓣宽\\
度 $x_1$ & 度 $x_2$ & 度 $x_3$ & 度 $x_4$ & 度 $x_1$ & 度 $x_2$ & 度 $x_3$ & 度 $x_4$ & 度 $x_1$ & 度 $x_2$ & 度 $x_3$ & 度 $x_4$\\
\midrule
5.1 & 3.5 & 1.4 & 0.2 & 7.0 & 3.2 & 4.7 & 1.4 & 6.3 & 3.3 & 6.0 & 2.5\\
4.9 & 3.0 & 1.4 & 0.2 & 6.4 & 3.2 & 4.5 & 1.5 & 5.8 & 2.7 & 5.1 & 1.9\\
4.7 & 3.2 & 1.3 & 0.2 & 6.9 & 3.1 & 4.9 & 1.5 & 7.1 & 3.0 & 5.9 & 2.1\\
4.6 & 3.1 & 1.5 & 0.2 & 5.5 & 2.3 & 4.0 & 1.3 & 6.3 & 2.9 & 5.6 & 1.8\\
5.0 & 3.6 & 1.4 & 0.2 & 6.5 & 2.8 & 4.6 & 1.5 & 6.5 & 3.0 & 5.8 & 2.2\\
\bottomrule
\end{tabular}
}
\end{mvtable}
